
\newpage

\stepcounter{section}
\section*{Lecture 4: Sheaf of Rings and Localisation}

In this lecture (number 4, originally given 6 August 2026), we define the Sheaf of Rings (and some elementary properties), and define and study the beginnings of Localisation.

Recall Grothenstieck: The topology carries the rings.

The Sheaf of Rings on the topological space of an affine scheme is called \emph{The Structure Sheaf}.
We recall from Lecture 3 that we need to associate a ring for every open set in $\Spec{R}$.
We showed that the empty set must be sent to a singleton, which for a ring means the zero ring.
We also recall that $\Spec{R}$ is the largest open set in $\Spec{R}$\footnote{
    Recall that the open sets are the complement of closed sets, so an open set is of the form $\Spec{R}\setminus{V(I)}$. Let $I=(1)$, then $V(I)$ is empty, so $\Spec{R}$ is the largest open subset.}, 
    so we can safely associate $\Spec{R}\rightsquigarrow R$.

This is supposed to be the analogue of smooth or holomorphic functions, so how do we consider the elements of $R$ as functions?

\begin{example}
    Let $R=\C[x,y]$, and note that these are polynomials over $x$ and $y$.
    We can consider the notion of evaluation of $f\in R$ by subsititution: the \emph{Value of $f$ at $p$} is $f(a,b)\in\C$ (where $f$ is a $\C$-valued function).
    This doesn't generalise well: how do you evaluate $67\in\Z$ (or any other member of $\Z$) in $\Spec\Z$? There is no polynomial for subsititution or evaluation.
\end{example}

Consider another approach that generalises well.
We recall that the underlying set (`points') of $\Spec{R}$ is the prime ideals of $\C[x,y]$ (consisting of $(0)$, ideals of irreducible polynomials in $\C[x,y]$ (see Hartshorne pg 7, Krull's \emph{Hauptidealsatz}, note that $\C[x,y]$ is a UFD), and the maximal ideals $(x-a,y-b)$).
For the maximal ideals, we have that $\mathfrak{m}\leftrightarrow \A_\C^2$.
We now consider the image of $f\in\C[x,y]$ in $\faktor{\C[x,y]}{(x-a,y-b)}$.
As $(x-a,y-b)$ is a maximal ideal, this gives us a field, and there's a canonical isomorphism of $\C$-algebras to this ring given by $[f]\mapsto f(a,b)$.
In this ideal, we note that $x\mapsto a$ and $y\mapsto b$ gives us what we want.
This is how we can think of elements of $R$ as functions.
\begin{example}
    Let $R=\C[x,y]$ and let $f\in\C[x,y]$.
    Then the evaluation at the maximal ideal $(x-a,y-b)$ corresponding to $(a,b)\in\A^2$ of $\Spec{R}$ is the image of $f$ in $\faktor{\C[x,y]}{(x-a,y-b)}$.
    As an example, consider $f=x-y$ at the points $(x-1,y-1)$ and $(x-1,y-2)$.
    Then at $(x-1,y-1)$, $f=(1-1)=[0]$ and at $(x-1,y-2)$ $f=(1-2)=[-1]$, which agrees with our intuition for evaluation.
\end{example}

\begin{example}
    Let $R=\Z$ and let $f=67$.
    Then evaluating $67$ at $(2)$ gives $[1]$.
    The \emph{Vanishing Locus} of $V(67)$ is the set of all primes at which $67$ has value $[0]$, which is only $67$: $V(67)=(67)$.
    $V(100)=\{(2),(5)\}$.
\end{example}

\begin{remark}
    $\Spec{\Z}$ is a really funky place to do geometry: $\Spec{\Z}$ describes a non-singular curve (Vakil).
\end{remark}

The traditional examples of these rings are things like:
\begin{align*}
    &\faktor{\C[x,y]}{(y^2-x^3)}\\
    &\faktor{\C[x,y,z]}{xy-z,zx-1}
\end{align*}

From the Nullstellensatz, we have that the maximal ideals of $\faktor{\C[x,y]}{(y^2-x^3)}$ are exactly those $(x-a,y-b)\ni (y^2-x^3)$.
These points of the scheme are closed, and are the \emph{closed points of the scheme}.
\begin{proof}
    Recall that the prime ideals are the points of $\Spec{R}$ and that a maximal ideal is prime.
    Then closed sets are `indexed' by the ideals: $V(\mathfrak{m})=\{\mathfrak{p}\mid \mathfrak{m}\subseteq\mathfrak{p}\}$.
    Clearly, we are looking for those ideals $P$ such that $\mathfrak{m}\subseteq P\subseteq\Spec{R}$.
    Note that $\mathfrak{m}$ is explicitly maximal, so $P$ is either $\mathfrak{m}$ or $R$.
    But $R$ is not considered a prime ideal by our convention, so $V(\mathfrak{m})=\{\mathfrak{m}\}$.
\end{proof}

If $\mathfrak{m}$ is maximal, $\faktor{R}{\mathfrak{m}}$ is the residue field at a closed point of $\Spec{R}$.
Quotienting by a maximal ideal of $\Spec{R}$ at a closed point gives us the base field of the scheme (over an algebraically closed field).

\begin{example}
    Let $R=\faktor{\Z}{24\Z}$. Take $f=12$.
    The points of $\Spec{R}$ are $(2)$ and $(3)$ because we have zero divisors (so $(0)$ is not prime).
    $f(2)=f(3)=0$. In fact, $f$ evaluated at any point of the scheme is zero, but $f\neq 0$ (is distinct from the zero function).
\end{example}

This is a genuinely new feature, and can happen in a more traditional example.

\begin{example}
    Let $R=\faktor{\C[x]}{x^3}$ and $f=x$.
    What is the value of $f$ at the point $(x)$ (a maximal ideal of $\faktor{\C[x]}{x^3})$?
    We want to see what the image of $f$ is in $\faktor{R}{(x)}$ (which is clearly $[0]$).
    So evaluating $f$ at a point of $\Spec{R}$ is a lossy thing.
\end{example}

At this point, we have assigned the ring $R$ to the largest open subset of $\Spec{R}$.

\subsection{Localisation}
    We need some new algebra to continue.

    Take a ring $R$, and a subset $S\subset R$ which is closed under multiplication.
    We will return to the examples of $S=\{1,f,f^2,f^3,\dots\}$ and $S$ the complement of a prime ideal (for example, $R=\Z$ and $S=(2)$, and $R_{(S)}$ is thus the odd numbers).
    Note that $S$ will always need to contain the unit; a short (but not satisfying) argument is that the product of zero elements must be $1$.
    
    We form a new ring: $S^{-1}R$ formed by equivalence classes of $(r_1,s_1)\sim(r_2,s_2)$ if there exists $s\in S$ such that $s\cdot(r_1s_2-s_1r_2)=0$.
    Considering $S^{-1}R$ as the fractions with denominators coming from $S$ is fine when $R$ is an integral domain, but falls apart if it is not (although it is a useful notation) via the equivalence relation.
    Using the previous definition fixes this issue.

    \begin{example}[How $S^{-1}R$ breaks when $R$ not a domain]
        We consider the equivalence relation defined by $(r_1,s_1)\sim(r_2,s_2)$ if $r_1s_2-s_1r_2=0$.
        Assume that $R$ is not an integral domain. Then there must exist zero divisors.
        We have that $(r_1,s_1)\sim(r_2,s_2)\implies r_1s_2=s_1r_2$, and $(r_2,s_2)\sim(r_3,s_3)\implies r_2s_3=s_2r_3$.
        Then we want to show that $(r_1,s_1)\sim(r_3,s_3)$.
        So we multiply through by $s_3$: $s_3r_1s_2-s_3s_1r_2=0$ and $s_1r_2s_3-s_1s_2r_3=0$.
        Note that we now have $s_3r_1s_2-s_1s_2r_3=0$, but if $s_2$ is a zero divisor, then $r_1s_3-s_1r_3$ may not be zero but may appear that way.
    \end{example}

    We note that $R\to S^{-1}R$ as $S$ contains the unit, so $R\subseteq S^{-1}R$. This is a natural map.
    When does $a\neq0\in R\mapsto0$ in $S^{-1}R$?
    This happens if and only if $(a,1)\sim(0,1)$ and this occurs if and only if $\exists sa=0$, which implies that $S$ has got zero divisors.
    This shows that the map is not always injective.

    More generally, if $S$ contains no zero divisors, then the map is injective.
    This map is very rarely surjective.
    Finally, the it can happen that the localisation is the zero ring, which occurs if and only if $0\in S$, which occurs if and only if $S$ contains a nilpotent.

    All the elements of $S$ are sent to units: the inverse of $(s,1)$ is $(1,s)$ (consider this written as a fraction).
    This is the smallest extension of $R$ such that all elements of $S$ become units.

    We now consider a universal mapping property:
    \begin{equation*}
        \begin{tikzcd}
            R\ar[d]\ar[dr,"\varphi"]& \\
            S^{-1}R\ar[r,dashed,"\exists!"]&B
        \end{tikzcd}
    \end{equation*}

    with $\varphi$ any ring homomorphism such that images of elements of $S$ are invertible in $B$.

    \subsection{Prime Ideals of Localisation}
        What are the prime ideals of $S^{-1}R$ and how does $\Spec{R}$ relate to $\Spec{S^{-1}R}$?
        \begin{equation*}
            \begin{tikzcd}
                R\ar[r,"i"]&S^{-1}R\\
                i^{-1}(\mathfrak{p})\ar[u, phantom, sloped, "\subset"]&\mathfrak{p}\ar[u, phantom, sloped, "\subset"]
            \end{tikzcd}
        \end{equation*}
        We first consider the map above, and note that $i^{-1}(\mathfrak{p})$ must be a prime ideal that does not contain an element of $S$, as elements of $S$ are sent to units in $S^{-1}R$ (by the correspondence theorem).
        Taking a prime ideal $\mathfrak{p}\in R$, $\l i(\mathfrak{p})\r=\left\{\frac{r}{s}\mid r\in\mathfrak{p}\right\}$.
        
        This is a bijection.
        Note that $\frac{f_1}{s_1}\cdot\frac{f_2}{s_2}=\frac{f}{s}$ for some $f$ and $s$.
        \begin{proof}
            By definition, $s^\prime sf_1f_2=s_1s_2fs^\prime$ (the equivalence relation), and is an equality in $R$.
            The RHS is a member of $\mathfrak{p}$, which gives us that the LHS must also be a member of $\mathfrak{p}$.
            But $s^\prime$ is not a member of $p$, which means that the other parts must be, so one of $f_1$ or $f_2$ must be in $\mathfrak{p}$
        \end{proof}

        What is $\Spec{R_f}$? (recall that $R_f=(1,f,f^2,\dots)^{-1}R$, the \emph{localisation of $R$ at $f$}).
        So we've got that $\Spec{R_f}=\Spec{R}\setminus{V(f)}\eqqcolon D_f$ (base of the topology), noting that $V(f)$ is the set of prime ideals where $f$ vanishes.
        So $\Spec{R_f}$ is the prime ideals of $R$ that don't contain $f$.

    \subsection{Building the Structure Sheaf}
        So we may now send $\Spec{R}\to R$ and $D_f\to R_f$ as the start of building the structure sheaf.
        If $D_f\subset D_g$, then we need to find a map $R_g\to R_f$.
        We need $(g)=(1)$ in $R_f$, otherwise we don't have a map.

        Take $\mathfrak{p}\subset R_f$ containing $g$ and look at the preimage $\tilde{\mathfrak{p}}\subset R$ of $\mathfrak{p}$.
        We note that $\tilde{\mathfrak{p}}$ contains $g$, but not $f$ as we're pulled back from a place where $f$ is invertible.

        We now have a pre-sheaf on the distinguished open sets, and this is the natural approach here.

        Recall the sheaf conditions:
        \begin{enumerate}
            \item If $f,g\in\mathcal{F}$ and $\{U_i\}$ is an open cover of some open $U$, then if $f\in\mathcal{F}(U_i)=g\in\mathcal{F}(U_i)$ for all $i$, then $f=g$
            \item Given $f_i\in\mathcal{F}(U_i)$ for all $i$ such that for all $i,j$ the image of $f_i\in\mathcal{F}(U_i\cap U_j)$ equals the image of $f_j\in\mathcal{F}(U_i\cap U_j)$ then there exists $f\in\mathcal{F}(U)$ that maps to all the $f_i$.
        \end{enumerate}

        If there is no prime ideal containig both $f$ and $g$, then $\Spec{R}=D_f\cup D_g$, then there are maps $R\to R_g$ and $R\to R_f$.
        The first condition requires for every ideal that $f\neq 0$ or $g\neq 0$, as there is no ideal with $f=0$ and $g=0$.
        Then $f\cup g=(1)$ and $(f,g)=(1)$, so we can find some $af+bg=1$.

        Take $h_1$ and $h_2$ such that they are equal on both localisations.
        So $(h_1-h_2)\cdot g^n=0$ and $(h_1-h_2)\cdot f^m=0$. Note that the differences are killed by powers of generators.
        We want $h_1-h_2=0$ (but not as a function that evaluates to zero, actually zero).
        For example, consider $n\cdot 2^m=0$ implies that $n=0$ (we don't necessarily have no nilpotents, so this is an important check).

        Two sections that are equal locally must be equal globally.
        We can check that the object that we have so far satisfies the sheaf axioms (Exercise.)

        We now have the structure sheaf:
        \begin{itemize}
            \item $\Spec{R}$ with this structure sheaf is our affine scheme
            \item Any open set can be covered by $D_f$s as they are a basis for the topology
            \item We need to check on double intersections
            \item This is effectively a book-keeping exercise
        \end{itemize}
